\section{Conclusion}

We presented \ourmethod{}, a text-space optimizer that treats an external skill document as the trainable state for frozen LLM agents. By separating the target model that executes tasks from the optimizer that edits skills, and by using bounded edit budgets, minibatch reflection, held-out validation gates, rejected-edit buffers, and epoch-wise slow/meta update, \ourmethod{} turns skill improvement into a controlled learning process rather than ad hoc prompt revision. Across six benchmarks, seven target models, and three execution modes, \ourmethod{} is best or tied-best on $52$ of $52$ evaluated cells, lifts GPT--5.5 by $+23.5$ points on average over no skill in direct chat and by $+24.8 / +19.1$ points under Codex and Claude Code harnesses, and beats the strongest per-cell baseline from human, LLM, Trace2Skill, TextGrad, GEPA, and EvoSkill skills by $+5.4$ points on average. Per-benchmark case studies show that these gains arise from compact ($<2{,}000$ token), interpretable skill artifacts assembled from only $1$--$4$ accepted edits, and that the deployed skills transfer across model scales, harnesses, and nearby benchmarks. These results suggest that compact natural-language skills can serve as a practical domain-adaptation layer for frontier agents, enabling reusable improvement without modifying model weights.

\paragraph{Outlook.}
\ourmethod{} optimizes a single skill artifact for a single target domain; natural extensions include skill libraries that share infrastructure across domains, reuse of optimizer-side meta skills across benchmarks, reward-free or preference-driven validation gates for open-ended tasks, and self-distillation of optimized skills back into the target model as a stepping stone toward weight-level adaptation. We hope that treating the skill itself as the trainable object---rather than as a side artifact of prompting---will let future work apply the full toolkit of optimization (learning rates, schedules, regularization, curricula, validation) to a part of the agent stack that has so far been hand-engineered.
